\section{What Your PA Expects}
\textit{Reference $|$ Course Text: Chapters 11, 13, 17}

Your PA evaluates you based on both \emph{what} you deliver and \emph{how} you deliver it. Meeting the checklist is necessary but not sufficient.

\begin{faqbox}[What separates an A team from a B team?]
\textbf{A teams}: communicate proactively (PA hears about problems before they become crises), incorporate feedback aggressively (PDR weaknesses are addressed by CDR), demonstrate engineering judgment (trade-offs are analyzed with data, not guesses), and produce professional-quality documentation (reports read like they were written by engineers, not by students rushing at midnight).

\textbf{B teams}: meet deadlines but do not exceed expectations, address feedback partially, make design decisions without data, and produce adequate but unpolished documentation.
\end{faqbox}

\begin{protip}
The single most effective way to improve your PA's assessment: \textbf{close the feedback loop}. When your PA gives feedback at a sprint review, write it down, address it, and report back at the next meeting: ``At the last review you asked about our thermal analysis. Here are the results.'' PAs notice when their feedback is taken seriously.
\end{protip}


% ═══════════════════════════════════════
